\chapter{Results}

Summarize the findings without interpretation. Results should be clear and concise. Only analyzed data should be presented in forms of figures, graphs, tables and/or text descriptions of observations. Avoid using phrases that are vague or non-specific, such as, "appeared to be greater than other variables..." or "demonstrates promising trends that...." Subjective modifiers should be explained in the discussion section of the paper [i.e., why did one variable appear greater? Or, how does the finding demonstrate a promising trend?].

When presenting statistically summarized data, you should state whether the number is a mean or median and clearly state how the data spread is expressed (± standard deviation, ± standard error of the mean, or inter-quartile range). When claiming a statistically significant result, you must support such a statement with a declaration of the probability (p) value and the test that was used to generate that value. 

When presenting results from your study, use the past tense when referring to your results.

All Figures and Tables should be numbered chronologically as they appear in your thesis. All Figures and Tables must be referred to in the text to facilitate reading. And they should be displayed in embedded data format, not compressed photos/ pictures. See further guidelines for constructing tables and figures in Part 3. 

Use the past tense when referring to your results.


\section{Evaluation Setup}

% Your evaluation setup description goes here.

\section{Main Results}

\subsection{Result Category 1}

% Your results content goes here.

% Example figure (uncomment and add your figure file):
% \begin{figure}[htbp]
%     \centering
%     \includegraphics[width=14.5cm]{figures/your_figure.pdf}
%     \caption{Your figure caption goes here.}
%     \label{fig:example}
% \end{figure}

\subsection{Result Category 2}

% Your results content goes here.

\section{Quantitative Analysis}

\subsection{Statistical Analysis}

% Your statistical analysis goes here.

\begin{table}[htbp]
    \centering
    \begin{tabular}{ll}
    \toprule
    Column 1 & Column 2 \\
    \midrule
    Data 1 & Data 2 \\
    \bottomrule
    \end{tabular}
    \caption{Your table caption goes here.}
    \label{tab:example}
\end{table}

\section{Summary of Observed Results}

% Your summary goes here.

